\hypertarget{c23-latest-features}{%
\chapter{C++23 LATEST FEATURES}\label{c23-latest-features}}

\hypertarget{c23-overview-direction}{%
\section{C++23 Overview \& Direction}\label{c23-overview-direction}}

C++23 (finalized in 2023) is a \textbf{refinement and enhancement} of
C++20 with practical improvements.

\hypertarget{timeline-context}{%
\subsection{Timeline \& Context}\label{timeline-context}}

\begin{itemize}
\tightlist
\item
  \textbf{2011}: C++11 (revolutionary)
\item
  \textbf{2014}: C++14 (refinement)
\item
  \textbf{2017}: C++17 (major improvements)
\item
  \textbf{2020}: C++20 (revolutionary leap)
\item
  \textbf{2023}: C++23 (practical enhancements)
\end{itemize}

\hypertarget{c23-philosophy}{%
\subsection{C++23 Philosophy}\label{c23-philosophy}}

\begin{itemize}
\tightlist
\item
  \textbf{Enhance} existing C++20 features
\item
  \textbf{Fill gaps} in C++20 design
\item
  \textbf{Improve} convenience and usability
\item
  \textbf{Optimize} common patterns
\item
  \textbf{Standardize} frequently-requested features
\item
  \textbf{Fix} issues discovered in C++20
\end{itemize}

\hypertarget{key-themes}{%
\subsection{Key Themes}\label{key-themes}}

\begin{enumerate}
\def\labelenumi{\arabic{enumi}.}
\tightlist
\item
  \textbf{Output \& Formatting} - std::print for easy output
\item
  \textbf{Error Handling} - std::expected for results
\item
  \textbf{Loop Control} - Enhanced for loops with ranges
\item
  \textbf{Memory Safety} - Better pointer/array handling
\item
  \textbf{Debugging} - Stack traces
\item
  \textbf{Templates} - Deducing this improvements
\item
  \textbf{Constexpr} - More compile-time power
\item
  \textbf{Library} - Quality of life improvements
\end{enumerate}

\hypertarget{why-c23-matters}{%
\subsection{Why C++23 Matters}\label{why-c23-matters}}

C++23 builds on C++20 strengths: - Easier output without iostream
overhead - Type-safe error handling (std::expected) - Better for loop
control - Debugging support (stack traces) - More flexible subscript
operator - Improved constexpr capabilities - More convenient library
features - Better optional support

\begin{center}\rule{0.5\linewidth}{0.5pt}\end{center}

\hypertarget{stdprint-formatted-output}{%
\section{STD::PRINT \& FORMATTED
OUTPUT}\label{stdprint-formatted-output}}

\hypertarget{stdprint---simple-output}{%
\section{1.1 std::print - Simple
Output}\label{stdprint---simple-output}}

\passthrough{\lstinline!std::print!} provides easy, fast output without
iostream overhead.

\hypertarget{basic-print-usage}{%
\subsection{Basic print Usage}\label{basic-print-usage}}

\begin{lstlisting}[language={C++}]
#include <print>
#include <iostream>

// Simple output (no newline by default)
std::print("Hello, World!");

// With newline
std::println("Hello, World!");

// With format
std::println("Number: {}, Float: {:.2f}", 42, 3.14159);

// To stderr
std::print(std::cerr, "Error: {}\n", "something went wrong");
std::println(std::cerr, "Error: {}", "something went wrong");
\end{lstlisting}

\hypertarget{print-vs-format-vs-iostream}{%
\subsection{print vs format vs
iostream}\label{print-vs-format-vs-iostream}}

\begin{lstlisting}[language={C++}]
#include <print>
#include <format>
#include <iostream>

std::string msg = "Hello";
int value = 42;

// iostream (slow, verbose)
std::cout << msg << ": " << value << "\n";

// format (creates string, then print)
std::cout << std::format("{}: {}\n", msg, value);

// print (direct output, fast)
std::println("{}: {}", msg, value);
\end{lstlisting}

\hypertarget{print-with-file-streams}{%
\subsection{print with File Streams}\label{print-with-file-streams}}

\begin{lstlisting}[language={C++}]
#include <print>
#include <fstream>

std::ofstream file("output.txt");

// Direct to file
std::println(file, "Line 1: {}", 42);
std::println(file, "Line 2: {}", "test");

// Simpler than:
// file << "Line 1: " << 42 << "\n";
\end{lstlisting}

\begin{center}\rule{0.5\linewidth}{0.5pt}\end{center}

\hypertarget{stdformat-enhancements}{%
\section{1.2 std::format Enhancements}\label{stdformat-enhancements}}

\hypertarget{format-improvements-in-c23}{%
\subsection{Format Improvements in
C++23}\label{format-improvements-in-c23}}

\begin{lstlisting}[language={C++}]
#include <format>

// More format options
double pi = 3.14159;

std::format("{:.2%}", 0.25);         // "25.00%" (percentage)
std::format("{:g}", 0.0001);         // General format
std::format("{:#x}", 255);           // "0xff" (with prefix)
std::format("{:_^10}", "test");      // "_____test" (custom fill)
\end{lstlisting}

\begin{center}\rule{0.5\linewidth}{0.5pt}\end{center}

\hypertarget{deducing-this}{%
\section{DEDUCING THIS}\label{deducing-this}}

\hypertarget{explicit-member-function-parameters}{%
\section{2.1 Explicit Member Function
Parameters}\label{explicit-member-function-parameters}}

\passthrough{\lstinline!Deducing this!} allows capturing the type and
constness of the object.

\hypertarget{basic-deducing-this}{%
\subsection{Basic Deducing This}\label{basic-deducing-this}}

\begin{lstlisting}[language={C++}]
#include <iostream>
using namespace std;

struct Counter {
    int count = 0;
    
    // Traditional
    void increment() {
        count++;
    }
    
    // With deducing this
    void increment_new(this auto& self) {
        self.count++;
    }
    
    // Value vs reference overloads now simple
    auto get_data(this auto& self) {
        return self.count;
    }
};

Counter c;
c.increment_new();
cout << c.get_data() << "\n";  // 1
\end{lstlisting}

\hypertarget{deducing-this-for-constnon-const}{%
\subsection{Deducing This for
Const/Non-Const}\label{deducing-this-for-constnon-const}}

\begin{lstlisting}[language={C++}]
struct Data {
    int value = 42;
    
    // Single function handles both const and non-const
    auto& get(this auto& self) {
        return self.value;
    }
    
    // Before C++23: Need two overloads
    // int& get() { return value; }
    // const int& get() const { return value; }
};

Data d;
d.get() = 100;              // Mutable reference
cout << d.get() << "\n";    // 100

const Data cd;
cout << cd.get() << "\n";   // 100 (const reference)
\end{lstlisting}

\hypertarget{practical-deducing-this}{%
\subsection{Practical Deducing This}\label{practical-deducing-this}}

\begin{lstlisting}[language={C++}]
template<typename T>
struct Optional {
    T value;
    bool has_value_flag = false;
    
    // Works for both optional<T> and optional<const T>
    auto* get_if_value(this auto* self) {
        return self->has_value_flag ? &self->value : nullptr;
    }
};

Optional<int> opt;
opt.value = 42;
opt.has_value_flag = true;

if (auto* ptr = opt.get_if_value()) {
    cout << *ptr << "\n";  // 42
}
\end{lstlisting}

\hypertarget{deducing-this---beyond-the-basics}{%
\section{2.2 Deducing This - Beyond the
Basics}\label{deducing-this---beyond-the-basics}}

\hypertarget{recursive-lambdas}{%
\subsection{Recursive Lambdas}\label{recursive-lambdas}}

Previously, lambdas couldn't easily call themselves. Now they can via
the explicit object parameter.

\begin{lstlisting}[language={C++}]
auto fib = [](this auto&& self, int n) {
    if (n <= 1) return n;
    return self(n - 1) + self(n - 2);
};

cout << fib(10) << "\n"; // 55
\end{lstlisting}

\hypertarget{replacing-crtp}{%
\subsection{Replacing CRTP}\label{replacing-crtp}}

The Curiously Recurring Template Pattern (CRTP) was used to inject
functionality into derived classes.
\passthrough{\lstinline!Deducing this!} simplifies it.

\textbf{Old CRTP:}

\begin{lstlisting}[language={C++}]
template <typename Derived>
struct Addable {
    Derived& operator+=(const Derived& other) {
        static_cast<Derived*>(this)->value += other.value;
        return *static_cast<Derived*>(this);
    }
};
struct Int : Addable<Int> { int value; };
\end{lstlisting}

\textbf{New C++23 Way:}

\begin{lstlisting}[language={C++}]
struct Addable {
    template <typename Self>
    auto& operator+=(this Self&& self, const Self& other) {
        self.value += other.value;
        return self;
    }
};
struct Int : Addable { int value; }; // No template parameter needed!
\end{lstlisting}

\begin{center}\rule{0.5\linewidth}{0.5pt}\end{center}

\hypertarget{range-based-for-loop-enhancements}{%
\section{RANGE-BASED FOR LOOP
ENHANCEMENTS}\label{range-based-for-loop-enhancements}}

\hypertarget{for-loop-initializers}{%
\section{3.1 For Loop Initializers}\label{for-loop-initializers}}

C++23 allows initialization in range-based for loops.

\hypertarget{for-loop-with-init}{%
\subsection{For Loop with Init}\label{for-loop-with-init}}

\begin{lstlisting}[language={C++}]
#include <vector>
#include <iostream>
using namespace std;

vector<int> v1 = {1, 2, 3};
vector<int> v2 = {4, 5, 6};

// Traditional
vector<int> combined;
for (int x : v1) combined.push_back(x);
for (int x : v2) combined.push_back(x);

// C++23: Initialize in loop
for (auto v = vector<int>{1, 2, 3, 4, 5, 6}; int x : v) {
    cout << x << " ";
}

// More practical
for (auto file = open_file("data.txt"); auto line : file.lines()) {
    cout << line << "\n";
}
\end{lstlisting}

\hypertarget{for-loop-with-init-and-structured-binding}{%
\subsection{For Loop with Init and Structured
Binding}\label{for-loop-with-init-and-structured-binding}}

\begin{lstlisting}[language={C++}]
map<string, vector<int>> data{
    {"a", {1, 2, 3}},
    {"b", {4, 5, 6}}
};

// Initialize and use with structured binding
for (auto it = data.begin(); auto [key, values] : data) {
    cout << key << ": ";
    for (int v : values) cout << v << " ";
    cout << "\n";
}
\end{lstlisting}

\begin{center}\rule{0.5\linewidth}{0.5pt}\end{center}

\hypertarget{stdexpected}{%
\section{STD::EXPECTED}\label{stdexpected}}

\hypertarget{result-type-for-error-handling}{%
\section{4.1 Result Type for Error
Handling}\label{result-type-for-error-handling}}

\passthrough{\lstinline!std::expected!} represents either a value or an
error.

\hypertarget{basic-expected-usage}{%
\subsection{Basic expected Usage}\label{basic-expected-usage}}

\begin{lstlisting}[language={C++}]
#include <expected>
#include <string>
#include <iostream>
using namespace std;

enum class ParseError { InvalidFormat, OutOfRange };

// Function returning expected
expected<int, ParseError> parse_int(const string& s) {
    try {
        return stoi(s);
    } catch (const invalid_argument&) {
        return unexpected(ParseError::InvalidFormat);
    } catch (const out_of_range&) {
        return unexpected(ParseError::OutOfRange);
    }
}

// Usage
auto result = parse_int("42");

if (result) {
    cout << "Value: " << result.value() << "\n";
} else {
    cout << "Error\n";
}
\end{lstlisting}

\hypertarget{expected-with-transform}{%
\subsection{expected with Transform}\label{expected-with-transform}}

\begin{lstlisting}[language={C++}]
#include <expected>

expected<int, string> get_value();

// Chain operations with transform
auto result = get_value()
    .transform([](int x) { return x * 2; })
    .transform_error([](const string& e) { return "Failed: " + e; });

if (result) {
    cout << result.value() << "\n";
} else {
    cout << result.error() << "\n";
}
\end{lstlisting}

\hypertarget{expected-vs-optional}{%
\subsection{expected vs optional}\label{expected-vs-optional}}

\begin{lstlisting}[language={C++}]
// optional: Has value or nothing
optional<int> opt = parse_int("42");
if (!opt) {
    // But why did it fail?
}

// expected: Has value or specific error
expected<int, ParseError> exp = parse_int("42");
if (!exp) {
    cout << "Error: " << static_cast<int>(exp.error()) << "\n";
}
\end{lstlisting}

\begin{center}\rule{0.5\linewidth}{0.5pt}\end{center}

\hypertarget{stdoptional-improvements}{%
\section{STD::OPTIONAL IMPROVEMENTS}\label{stdoptional-improvements}}

\hypertarget{enhanced-optional-operations}{%
\section{5.1 Enhanced optional
Operations}\label{enhanced-optional-operations}}

\hypertarget{optional-with-deref-operator}{%
\subsection{optional with Deref
Operator}\label{optional-with-deref-operator}}

\begin{lstlisting}[language={C++}]
#include <optional>

optional<int> opt{42};

// C++23: Monadic operations
auto result = opt
    .and_then([](int x) -> optional<int> { return x * 2; })
    .or_else([]() { return optional<int>(0); });

// C++23: Chaining
opt.transform([](int x) { return x + 10; })
   .and_then([](int x) -> optional<int> {
       if (x > 50) return x;
       return nullopt;
   });
\end{lstlisting}

\hypertarget{optionalvalue_or_else}{%
\subsection{optional::value\_or\_else}\label{optionalvalue_or_else}}

\begin{lstlisting}[language={C++}]
#include <optional>

optional<int> opt;

// Get value or call function to generate default
int value = opt.value_or_else([]() { return compute_default(); });

// More flexible than value_or
// value_or: int value = opt.value_or(0);
// value_or_else: int value = opt.value_or_else([]() { return expensive_computation(); });
\end{lstlisting}

\begin{center}\rule{0.5\linewidth}{0.5pt}\end{center}

\hypertarget{multidimensional-subscript-operator}{%
\section{MULTIDIMENSIONAL SUBSCRIPT
OPERATOR}\label{multidimensional-subscript-operator}}

\hypertarget{multiple-index-support}{%
\section{6.1 Multiple Index Support}\label{multiple-index-support}}

C++23 allows multiple indices in subscript operator.

\hypertarget{multi-index-subscript}{%
\subsection{Multi-Index Subscript}\label{multi-index-subscript}}

\begin{lstlisting}[language={C++}]
#include <iostream>
using namespace std;

struct Matrix {
    int data[10][10];
    
    // C++23: Multi-index operator
    int& operator[](int row, int col) {
        return data[row][col];
    }
    
    int& operator[](int row, int col) const {
        return data[row][col];
    }
};

Matrix m;
m[2, 3] = 42;           // Two indices
cout << m[2, 3] << "\n"; // 42
\end{lstlisting}

\hypertarget{dynamic-2d-array-wrapper}{%
\subsection{Dynamic 2D Array Wrapper}\label{dynamic-2d-array-wrapper}}

\begin{lstlisting}[language={C++}]
template<typename T>
struct Array2D {
    vector<T> data;
    size_t width, height;
    
    Array2D(size_t w, size_t h) : width(w), height(h), data(w * h) {}
    
    T& operator[](size_t row, size_t col) {
        return data[row * width + col];
    }
    
    const T& operator[](size_t row, size_t col) const {
        return data[row * width + col];
    }
};

Array2D<int> grid(3, 3);
grid[1, 1] = 5;
cout << grid[1, 1] << "\n";  // 5
\end{lstlisting}

\begin{center}\rule{0.5\linewidth}{0.5pt}\end{center}

\hypertarget{stdmdspan-multidimensional-view}{%
\section{6.2 std::mdspan (Multidimensional
View)}\label{stdmdspan-multidimensional-view}}

\passthrough{\lstinline!std::mdspan!} provides a non-owning
multidimensional view of contiguous data.

\hypertarget{basic-mdspan-usage}{%
\subsection{Basic mdspan Usage}\label{basic-mdspan-usage}}

\begin{lstlisting}[language={C++}]
#include <mdspan>
#include <vector>
#include <iostream>

int main() {
    std::vector<int> data = {
        1, 2, 3, 4,
        5, 6, 7, 8,
        9, 10, 11, 12
    };

    // Create a 3x4 view over the data (C++23)
    // using MatrixView = std::mdspan<int, std::dextents<size_t, 2>>;
    // MatrixView m(data.data(), 3, 4);
    
    // m[1, 2] == 7 (row 1, col 2)
    // No copying of data involved!
    
    return 0;
}
\end{lstlisting}

\hypertarget{mdspan-layouts}{%
\section{6.3 mdspan Layouts}\label{mdspan-layouts}}

You can control how 2D indices map to the 1D memory.

\begin{itemize}
\tightlist
\item
  \passthrough{\lstinline!std::layout\_right!}: Row-major (default in
  C++). Index \passthrough{\lstinline!(i, j)!} is
  \passthrough{\lstinline!i * N + j!}.
\item
  \passthrough{\lstinline!std::layout\_left!}: Column-major
  (Fortran/MATLAB). Index \passthrough{\lstinline!(i, j)!} is
  \passthrough{\lstinline!j * M + i!}.
\end{itemize}

\begin{lstlisting}[language={C++}]
using RowMajor = std::mdspan<double, std::dextents<size_t, 2>, std::layout_right>;
using ColMajor = std::mdspan<double, std::dextents<size_t, 2>, std::layout_left>;

// Same data, different interpretation
std::vector<double> v = {1, 2, 3, 4}; // 2x2 matrix
RowMajor m_row(v.data(), 2, 2); // [[1, 2], [3, 4]]
ColMajor m_col(v.data(), 2, 2); // [[1, 3], [2, 4]]
\end{lstlisting}

\begin{center}\rule{0.5\linewidth}{0.5pt}\end{center}

\hypertarget{stdstacktrace}{%
\section{STD::STACKTRACE}\label{stdstacktrace}}

\hypertarget{runtime-stack-trace-capture}{%
\section{7.1 Runtime Stack Trace
Capture}\label{runtime-stack-trace-capture}}

\passthrough{\lstinline!std::stacktrace!} provides runtime stack trace
information.

\hypertarget{basic-stacktrace}{%
\subsection{Basic Stacktrace}\label{basic-stacktrace}}

\begin{lstlisting}[language={C++}]
#include <stacktrace>
#include <print>
#include <iostream>

void deep_function() {
    // Capture current stack trace
    auto trace = std::stacktrace::current();
    
    // Print trace
    std::println("Stack trace:");
    for (const auto& entry : trace) {
        std::println("  {}", entry);
    }
}

void middle_function() {
    deep_function();
}

void top_function() {
    middle_function();
}

int main() {
    top_function();
    return 0;
}

// Output:
// Stack trace:
//   top_function()
//   middle_function()
//   deep_function()
\end{lstlisting}

\hypertarget{stacktrace-in-error-handling}{%
\subsection{Stacktrace in Error
Handling}\label{stacktrace-in-error-handling}}

\begin{lstlisting}[language={C++}]
#include <stacktrace>
#include <exception>

class TracedException : public std::exception {
    std::stacktrace trace;
    std::string message;
    
public:
    TracedException(const std::string& msg) 
        : trace(std::stacktrace::current()), message(msg) {}
    
    const char* what() const noexcept override {
        return message.c_str();
    }
    
    const std::stacktrace& get_trace() const {
        return trace;
    }
};

// Usage
void operation() {
    if (error_condition) {
        throw TracedException("Operation failed");
    }
}

try {
    operation();
} catch (const TracedException& e) {
    std::println("Error: {}", e.what());
    std::println("Trace: {}", e.get_trace());
}
\end{lstlisting}

\begin{center}\rule{0.5\linewidth}{0.5pt}\end{center}

\hypertarget{constexpr-enhancements}{%
\section{CONSTEXPR ENHANCEMENTS}\label{constexpr-enhancements}}

\hypertarget{more-compile-time-capabilities}{%
\section{8.1 More Compile-Time
Capabilities}\label{more-compile-time-capabilities}}

\hypertarget{constexpr-stdstring}{%
\subsection{constexpr std::string}\label{constexpr-stdstring}}

\begin{lstlisting}[language={C++}]
#include <string>

// C++23: Can use std::string in constexpr context
constexpr std::string concat(const char* a, const char* b) {
    std::string result = a;
    result += b;
    return result;
}

constexpr auto msg = concat("Hello", " World");
// msg = "Hello World" at compile-time
\end{lstlisting}

\hypertarget{constexpr-vector-operations}{%
\subsection{constexpr vector
Operations}\label{constexpr-vector-operations}}

\begin{lstlisting}[language={C++}]
#include <vector>

// C++23: vector works in constexpr
constexpr std::vector<int> make_sequence(int n) {
    std::vector<int> result;
    for (int i = 0; i < n; i++) {
        result.push_back(i);
    }
    return result;
}

constexpr auto seq = make_sequence(5);
// seq = {0, 1, 2, 3, 4} at compile-time
\end{lstlisting}

\hypertarget{constexpr-algorithms}{%
\subsection{constexpr Algorithms}\label{constexpr-algorithms}}

\begin{lstlisting}[language={C++}]
#include <algorithm>

// C++23: Algorithms work in constexpr
constexpr int compute() {
    std::vector<int> v{3, 1, 4, 1, 5, 9};
    std::sort(v.begin(), v.end());
    int sum = 0;
    for (int x : v) sum += x;
    return sum;
}

constexpr int result = compute();  // Computed at compile-time
\end{lstlisting}

\begin{center}\rule{0.5\linewidth}{0.5pt}\end{center}

\hypertarget{adaptor-improvements}{%
\section{ADAPTOR IMPROVEMENTS}\label{adaptor-improvements}}

\hypertarget{rangesto-conversion}{%
\section{9.1 Ranges::to Conversion}\label{rangesto-conversion}}

\passthrough{\lstinline!std::ranges::to!} converts ranges to containers.

\hypertarget{basic-rangesto}{%
\subsection{Basic ranges::to}\label{basic-rangesto}}

\begin{lstlisting}[language={C++}]
#include <ranges>
#include <vector>
#include <string>

vector<int> v = {1, 2, 3, 4, 5};

// Convert filtered range to vector
auto evens = v
    | std::ranges::views::filter([](int x) { return x % 2 == 0; })
    | std::ranges::to<std::vector>();

// Or with map
map<int, int> m{{1, 10}, {2, 20}, {3, 30}};
auto keys = m
    | std::ranges::views::keys
    | std::ranges::to<std::vector>();
\end{lstlisting}

\hypertarget{rangesto-with-construction-args}{%
\subsection{ranges::to with Construction
Args}\label{rangesto-with-construction-args}}

\begin{lstlisting}[language={C++}]
#include <ranges>
#include <set>

vector<int> v = {3, 1, 4, 1, 5, 9, 2, 6};

// Convert to sorted unique set
auto unique_sorted = v
    | std::ranges::to<std::set>();

// Or to deque
auto d = v | std::ranges::to<std::deque>();

// With custom construction
auto s = v
    | std::ranges::to<std::set>(std::greater{});  // Descending
\end{lstlisting}

\begin{center}\rule{0.5\linewidth}{0.5pt}\end{center}

\hypertarget{library-improvements}{%
\section{LIBRARY IMPROVEMENTS}\label{library-improvements}}

\hypertarget{utility-improvements}{%
\section{10.1 Utility Improvements}\label{utility-improvements}}

\hypertarget{stdout_ptr-stdinout_ptr}{%
\subsection{std::out\_ptr \&
std::inout\_ptr}\label{stdout_ptr-stdinout_ptr}}

\begin{lstlisting}[language={C++}]
#include <memory>

// For converting unique_ptr to C API
unique_ptr<int> ptr;

// Before C++23: Complicated
// int* tmp = ptr.release();
// legacy_function(&tmp);
// ptr.reset(tmp);

// C++23: Simple
legacy_c_function(std::out_ptr(ptr));

// For bidirectional
inout_ptr(ptr);
\end{lstlisting}

\hypertarget{stdmove_iterator-improvements}{%
\subsection{std::move\_iterator
Improvements}\label{stdmove_iterator-improvements}}

\begin{lstlisting}[language={C++}]
#include <iterator>
#include <algorithm>

vector<string> v = {"a", "b", "c"};

// C++23: Cleaner move semantics
auto result = v
    | std::views::transform([](auto& s) { return std::move(s); });
\end{lstlisting}

\hypertarget{bit-manipulation-improvements}{%
\subsection{Bit Manipulation
Improvements}\label{bit-manipulation-improvements}}

\begin{lstlisting}[language={C++}]
#include <bit>

unsigned int x = 5;  // 0b0101

// C++23 additions
std::byteswap(x);              // Byte swap
std::has_single_bit(x);        // Check if power of 2
std::bit_width(x);             // Bits needed
std::popcount(x);              // Count 1 bits
\end{lstlisting}

\begin{center}\rule{0.5\linewidth}{0.5pt}\end{center}

\hypertarget{attributes-features}{%
\section{ATTRIBUTES \& FEATURES}\label{attributes-features}}

\hypertarget{assume-attribute}{%
\section{11.1 {[}{[}assume{]}{]} Attribute}\label{assume-attribute}}

\passthrough{\lstinline![[assume]]!} allows providing hints to
optimizer.

\hypertarget{using-assume}{%
\subsection{Using assume}\label{using-assume}}

\begin{lstlisting}[language={C++}]
void process(int x) {
    // Tell compiler to assume x > 0 (for optimization)
    [[assume(x > 0)]];
    
    // Compiler can optimize based on this
    if (x < 0) {
        // This branch won't be taken
        std::cout << "Negative\n";
    }
}

// Useful for performance-critical code
int* find_first(int* arr, int size) {
    [[assume(size > 0)]];  // Array has elements
    
    for (int i = 0; i < size; i++) {
        if (arr[i] == target) return &arr[i];
    }
    return nullptr;
}
\end{lstlisting}

\hypertarget{stdcall-and-abi-attributes}{%
\section{11.2 {[}{[}stdcall{]}{]} and ABI
Attributes}\label{stdcall-and-abi-attributes}}

\begin{lstlisting}[language={C++}]
// Platform-specific calling conventions
#ifdef _WIN32
void __stdcall legacy_function() { }
[[gnu::stdcall]] void c_function();
#endif
\end{lstlisting}

\begin{center}\rule{0.5\linewidth}{0.5pt}\end{center}

\hypertarget{standard-library-additions}{%
\section{STANDARD LIBRARY ADDITIONS}\label{standard-library-additions}}

\hypertarget{container-utility-additions}{%
\section{12.1 Container \& Utility
Additions}\label{container-utility-additions}}

\hypertarget{stddebug_assert-conditional-assertion}{%
\subsection{std::debug\_assert (Conditional
Assertion)}\label{stddebug_assert-conditional-assertion}}

\begin{lstlisting}[language={C++}]
void function(int x) {
    // Debug assertion (disabled in release)
    _ASSERT(x > 0, "x must be positive");
    
    // Work with x
}
\end{lstlisting}

\hypertarget{stdrepeat_view}{%
\subsection{std::repeat\_view}\label{stdrepeat_view}}

\begin{lstlisting}[language={C++}]
#include <ranges>

// Repeat a value
auto repeated = std::views::repeat(42, 5);
for (int x : repeated) {
    cout << x << " ";  // 42 42 42 42 42
}
\end{lstlisting}

\hypertarget{stdstride_view}{%
\subsection{std::stride\_view}\label{stdstride_view}}

\begin{lstlisting}[language={C++}]
#include <ranges>

vector<int> v = {0, 1, 2, 3, 4, 5, 6, 7, 8, 9};

// Take every Nth element
auto every_other = v | std::views::stride(2);
for (int x : every_other) {
    cout << x << " ";  // 0 2 4 6 8
}
\end{lstlisting}

\hypertarget{stdchunk_by}{%
\subsection{std::chunk\_by}\label{stdchunk_by}}

\begin{lstlisting}[language={C++}]
#include <ranges>

vector<int> v = {1, 2, 2, 3, 3, 3, 4, 4};

// Group by predicate
auto chunks = v | std::views::chunk_by(std::equal_to{});
for (auto chunk : chunks) {
    cout << "[";
    for (int x : chunk) cout << x << " ";
    cout << "] ";
}
// Output: [1 ] [2 2 ] [3 3 3 ] [4 4 ]
\end{lstlisting}

\hypertarget{stdflat_map-and-stdflat_set}{%
\subsection{std::flat\_map and
std::flat\_set}\label{stdflat_map-and-stdflat_set}}

\passthrough{\lstinline!std::flat\_map!} is a container adaptor that
stores elements in sorted order in contiguous memory (like a vector).

\begin{lstlisting}[language={C++}]
#include <flat_map>
#include <iostream>
#include <string>
#include <vector>

int main() {
    // Stores keys and values in separate vectors
    // Cache-friendly, fast iteration, binary search
    std::flat_map<int, std::string> map;
    
    map[1] = "one";
    map[3] = "three";
    map[2] = "two";  // Inserted in correct sorted position
    
    for (const auto& [key, val] : map) {
        std::cout << key << ": " << val << "\n";
    }
    // Output: 1: one, 2: two, 3: three
    
    return 0;
}
\end{lstlisting}

\hypertarget{stdgenerator-synchronous-coroutine}{%
\subsection{std::generator (Synchronous
Coroutine)}\label{stdgenerator-synchronous-coroutine}}

\passthrough{\lstinline!std::generator!} is the standard coroutine
generator for synchronous sequences.

\begin{lstlisting}[language={C++}]
#include <generator>
#include <iostream>
#include <ranges>

std::generator<int> fib(int n) {
    int a = 0, b = 1;
    while (n-- > 0) {
        co_yield a;
        auto next = a + b;
        a = b;
        b = next;
    }
}

int main() {
    for (int x : fib(10)) {
        std::cout << x << " ";
    }
    // Output: 0 1 1 2 3 5 8 13 21 34
    
    // Composable with ranges
    auto gen = fib(100) | std::views::filter([](int x) { return x % 2 == 0; });
    
    return 0;
}
\end{lstlisting}

\hypertarget{generator-internals}{%
\subsection{12.3 Generator Internals}\label{generator-internals}}

\passthrough{\lstinline!std::generator!} is a coroutine that: 1.
\textbf{Suspends on yield}: \passthrough{\lstinline!co\_yield!} suspends
execution and returns value to caller. 2. \textbf{Promise Type}: Handles
the \passthrough{\lstinline!yield\_value!} call. 3. \textbf{Recursive}:
Can \passthrough{\lstinline!co\_yield!} another generator (unlike basic
coroutines).

\begin{lstlisting}[language={C++}]
// Pseudocode of how it works
struct promise_type {
    int current_value;
    std::suspend_always yield_value(int value) {
        current_value = value;
        return {}; // Suspend
    }
};
\end{lstlisting}

\begin{center}\rule{0.5\linewidth}{0.5pt}\end{center}

\hypertarget{c23-best-practices}{%
\section{C++23 BEST PRACTICES}\label{c23-best-practices}}

\hypertarget{whats-better-with-c23}{%
\section{What's Better with C++23}\label{whats-better-with-c23}}

\begin{lstlisting}[language={C++}]
// 1. Use std::print for simple output
std::println("Value: {}", value);

// 2. Use std::expected for error handling
expected<int, Error> result = operation();

// 3. Use deducing this for simpler overloads
auto get(this auto& self) { return self.value; }

// 4. Use range-based for with init
for (auto data = load_data(); auto item : data) { }

// 5. Use multi-index subscript
matrix[row, col] = value;

// 6. Use std::stacktrace for debugging
auto trace = std::stacktrace::current();

// 7. Use constexpr string/vector
constexpr auto msg = std::string("hello");

// 8. Use std::ranges::to for conversions
auto set_result = v | std::ranges::to<std::set>();

// 9. Use [[assume]] for optimization hints
[[assume(pointer != nullptr)]];

// 10. Use monadic operations on optional
opt.transform([](int x) { return x * 2; });
\end{lstlisting}

\begin{center}\rule{0.5\linewidth}{0.5pt}\end{center}
